\chapter{Introduction}

Almost every organisation now depends on networked systems for its daily work. Servers, workstations, cloud services and applications exchange data continuously, and each of these components can be targeted by attackers. Defending such an environment means watching several kinds of evidence at once: the traffic moving across the network, the files that arrive on endpoints, and the events recorded in system and application logs. Guidance on incident handling treats detection and analysis as a core phase of incident response, and points out that it depends on collecting and interpreting evidence from many different sources \cite{nist80061r3}.

In practice this evidence is often examined with separate tools. A network intrusion detection system raises alerts about suspicious flows, an antivirus or endpoint product reports suspicious files, and a log management platform stores events for later search. Commercial security information and event management (SIEM) and extended detection and response (XDR) products such as Splunk Enterprise Security, Microsoft Sentinel and CrowdStrike Falcon bring much of this information together \cite{splunkes,sentineldocs,crowdstrikefalcon}, and open-source platforms such as Wazuh offer comparable monitoring features \cite{wazuhdocs}. Even so, an analyst must still decide which alerts matter, recognise what kind of attack they represent and judge whether unfamiliar behaviour is harmful. When alerts are numerous, or when an attack does not match an existing rule, this work becomes slow and depends heavily on individual expertise.

Machine learning has been studied for many years as a way to support these decisions. Supervised classifiers can learn the patterns of known attack families from labelled data \cite{buczak2016survey}, while anomaly detection methods learn what normal behaviour looks like and flag deviations from it \cite{chandola2009anomaly}. Each approach has well-known weaknesses. Supervised models are usually accurate on the attacks they were trained on but cannot recognise genuinely new ones, and anomaly detectors tend to produce more false alarms and are harder to interpret \cite{sommer2010outside}. Results obtained on benchmark datasets also do not always transfer to real networks \cite{sommer2010outside,arp2022dos}. These observations suggest that no single model is sufficient, and that different detectors should be combined in a way that keeps their outputs understandable to a human analyst.

This project, the \textit{AI-Based Autonomous Cyber Defense System}, takes that view. It is being developed as a modular platform in which separate machine learning components examine different kinds of security evidence: network flows, malware feature vectors and deviations from normal traffic. Each component returns a structured result that can later be combined with the others. The long-term plan, set out in the project proposal, also includes log analysis, incident correlation, vulnerability and attack-path analysis, risk estimation, an AI investigation agent and a dashboard. The word ``autonomous'' in the title refers to automating analysis and the preparation of recommended responses. The system is designed to assist analysts, not to take enforcement actions on production systems by itself.

This mid-term report covers the first half of the project. By this stage, three analysis modules have been implemented and evaluated on public datasets: network threat classification, malware classification and anomaly detection. Log analysis is under development, and the remaining modules and the integration layer are planned. Throughout the report, implemented functionality is clearly separated from planned functionality. The rest of this chapter describes the purpose, scope, objectives, intended market and intended users of the system.
